\documentclass[11pt]{article} 

\usepackage{amsmath}

\usepackage[super]{nth}

\usepackage{epsfig}

\usepackage{graphics}

\usepackage{wasysym}

\usepackage{times}

\usepackage{natbib}

%\usepackage{amssymb}

\setlength{\topmargin}{0in}
\setlength{\textwidth}{6.25in}

\setlength{\oddsidemargin}{.125in}

% 28 lectures total

\usepackage{amscd}

\begin{document}
\pagenumbering{gobble}
\title{Quiz \hfill Name : \underline{\hspace{3cm}}}
\date{November 14}
\maketitle
\begin{enumerate}
\item[Example 3.20] Suppose a binary signal $S\in{-1,\ 1}$ with $P(S=1)=p$ is transmitted, but the signal is received with standard normally distributed noise $N$, so that we only observe $Y=N+S$.  Bayes Rule says the probability that $S=1$ given an observed value $y$ is

$$P(S=1|Y=y)=\frac{P(S=1) f_{Y|S}(y|1)}{f_Y(y)}=\frac{\frac{p}{\sqrt{2\pi}}e^{-(y-1)^2/2}}{\frac{p}{\sqrt{2\pi}}e^{-(y-1)^2/2}+\frac{1-p}{\sqrt{2\pi}}e^{-(y+1)^2/2}}=\frac{pe^y}{pe^y+(1-p)e^{-y}}$$


\begin{enumerate}
\item What happens to the probability that the S=1 if $y$ becomes very negative?
\bigskip
\bigskip
\bigskip

\item What is the probability that the signal is 1 if $y=0$?
\bigskip
\bigskip
\bigskip
\bigskip

\item Why was it valid that $f_Y(y)=\frac{p}{\sqrt{2\pi}}e^{-(y-1)^2/2}+\frac{1-p}{\sqrt{2\pi}}e^{-(y+1)^2/2}$ in the denominator?
\bigskip
\bigskip
\bigskip
\bigskip

\end{enumerate}



\item True or False?

\begin{enumerate}
\item Given two random variables X and Y, the joint PDF must be such that $$\int_{-\infty}^{\infty} f_{X,Y}(x,\ y)dx=1$$
\item A valid pdf must take on finite values.
\item The expectation of a product of 2 random variables is the product of their expectations.
\item The expectation of a sum of 2 random variables is the sum of their expectations.
\item If X and Y are independent, var($X+Y$)=var($X$)+var($Y$)
\end{enumerate}


\end{enumerate}
\end{document}